% !TeX encoding = UTF-8
% !TeX spellcheck = ca_ES
%
% P3CTeX-architecture.tex
% High-level architecture overview for the P3CTeX ecosystem.
%
% This document is informational only; it does not define any user API.
%
\documentclass[EN]{P3CTeX}

\begin{document}

\section*{Arquitectura de P3CTeX}

P3CTeX es divideix en tres capes principals, cadascuna amb una
responsabilitat ben delimitada:

\begin{description}
  \item[\texttt{P3CTeX.cls}] Classe de document principal. Gestiona:
    \begin{itemize}
      \item El comportament global del document (format, pàgina,
            tipografia bàsica).
      \item Les dades de context: estudiant, assignatura, activitat,
            idioma, etc.
      \item La interfície d'usuari de nivell document (\LaTeX2e): ordres
            com ara la configuració de la portada i el sumari.
    \end{itemize}
    La classe \emph{no} carrega biblioteques independents directament
    (\texttt{pxUML}, \texttt{pxPRP}, \dots); ho fa sempre a través del
    paquet d'agregació \texttt{pxCORE}.

  \item[\texttt{pxCORE.sty}] Paquet nucli i d'agregació. S'encarrega de:
    \begin{itemize}
      \item Centralitzar el \emph{bootstrap} de paquets externs
            compartits (gràfics, tipografies, enllaços, etc.).
      \item Definir la família de claus \texttt{pxCORE}, que permet
            activar o desactivar biblioteques independents
            (\texttt{UML}, \texttt{PRP}, \dots) i passar-los opcions.
      \item Carregar les biblioteques independents (\texttt{pxUML},
            \texttt{pxPRP}, \texttt{pxSRC}) segons les
            claus definides.
    \end{itemize}
    El paquet \texttt{pxCORE} no defineix macros d'usuari
    específiques de cap classe de document; simplement ofereix un
    punt únic d'integració.

  \item[Biblioteques independents (\texttt{pxPRP}, \texttt{pxUML}, \dots)]
    Cada biblioteca encapsula una funcionalitat concreta:
    \begin{itemize}
      \item \texttt{pxPRP}: gestió d'objectes lògics basats en
            \texttt{prop} amb llistes separades per comes.
      \item \texttt{pxUML}: diagrames UML basats en \texttt{pgf-umlcd}
            i \texttt{TikZ}, recolzats en \texttt{pxPRP}.
    \end{itemize}
    Les biblioteques:
    \begin{itemize}
      \item Defineixen la seva API pública només mitjançant macros
            \LaTeX2e als fitxers \texttt{*.sty}.
      \item Concentren la lògica d'implementació interna en fitxers
            \texttt{*.code.tex} basats en \texttt{expl3}.
      \item No depenen de detalls interns de \texttt{P3CTeX.cls};
            poden ser utilitzades, en la mesura del possible, de
            manera autònoma.
    \end{itemize}
\end{description}

\bigskip

En resum:
\begin{center}
  \texttt{P3CTeX.cls} $\rightarrow$ (configura) \texttt{pxCORE.sty}
  $\rightarrow$ (activa) biblioteques \texttt{px*}.
\end{center}

\end{document}

